\section{Требования к системе}

На основе проведённого анализа существующих платформ были сформулированы требования к
разрабатываемому комплексу информационных систем.
Требования разделены на функциональные и нефункциональные.

\subsection{Функциональные требования}

\subsubsection*{Пользователи}
Система должна обеспечивать регистрацию и аутентификацию пользователей по никнейму и паролю.
Пользователь идентифицируется по уникальному никнейму.

\subsubsection*{Группы пользователей}
Система должна поддерживать создание и управление группами пользователей.
Группа идентифицируется по уникальному числовому идентификатору.
Пользователи в группе могут иметь роли создателя, менеджера и участника,
которые определяют их права на совершение действий в группе.
Пользователи могут быть добавлены в группу с определённой ролью, которая впоследствии может быть изменена.
Пользователи с более высокими правами могут исключать пользователей с меньшими правами из группы,
а также покинуть группу самостоятельно.

\subsubsection*{Рейтинг пользователя}
Система должна реализовывать рейтинговую систему.
Пользователь получает начальный рейтинг при регистрации.
Рейтинг меняется при участии пользователя в рейтинговых дуэлях.

\subsubsection*{Рейтинговая дуэль}
Система должна поддерживать рейтинговый режим дуэли.
Пользователь может начать поиск рейтинговой дуэли,
и система должна подобрать соперника, близкого по рейтингу.
Рейтинг меняется в соответствии с результатом дуэли и текущими рейтингами участников.

\subsubsection*{Дружеская дуэль}
Система должна поддерживать дружеский режим дуэли.
Пользователь может отправить приглашение на дружескую дуэль другому пользователю.
Второй пользователь имеет возможность принять и отклонить приглашение.
Рейтинг после дружеской дуэли не меняется.

\subsubsection*{Конфигурация дуэли}
Система должна поддерживать управление конфигурацией дуэли.
Пользователь может создать и отредактировать конфигурацию дуэли.
В конфигурацию дуэли входят следующие параметры:
\begin{itemize}
\item длительность дуэли;
\item возможность смотреть код соперника во время дуэли;
\item количество задач в дуэли;
\item этапность задач в дуэли (все задачи доступны сразу или открываются по очереди);
\item сложность и темы каждой задачи в дуэли.
\end{itemize}
Имеется стандартная конфигурация дуэли, по которой проводятся рейтинговые дуэли.
Пользователь имеет возможность использовать для дружеской дуэли любую из своих конфигураций.

\subsubsection*{Дуэль в группе}
Система должна предоставлять возможность создания дуэли в группе между определёнными пользователями.
Пользователи с определёнными ролями имеют возможность создать дуэль со своей конфигурацией между любыми двумя участниками группы.
Дуэль в таком случае является нерейтинговой.

\subsubsection*{Турнир в группе}
Система должна предоставлять возможность создания турнира в группе.
Турнир имеет схему проведения (например, олимпийская система или групповая стадия).
Для всех дуэлей в турнире можно выбрать единую конфигурацию дуэли.
Можно выбрать пользователей группы, которые будут участвовать в турнире.

\subsubsection*{Участие в дуэли}
Интерфейс участия в дуэли поделён на две части.
В одной отображается редактор кода, в котором пользователь должен писать своё решение задачи.
Во второй отображается условие задачи и наборы открытых тестовых данных.
Участник дуэли имеет возможность отправить текущее решение на проверку и получить вердикт тестирования.
В дуэли побеждает пользователь, который раньше другого отправит первое правильное решение.
Если по истечении длительности дуэли правильных решений отправлено не было, объявляется ничья.

\subsubsection*{Запуск решения на тестах}
Интерфейс участия в дуэли предоставляет возможность запустить своё решение на открытых тестах и увидеть выходные данные.
Также пользователь имеет возможность запустить своё решение на своём произвольном тесте.

\subsubsection*{Автоматическое тестирование решений}
Система должна автоматически тестировать решения пользователей.
Это подразумевает компиляцию и/или запуск пользовательского кода с последующей проверкой правильности ответа.
Система должна поддерживать выполнение решений на нескольких языках программирования.
Проверка корректности ответа может осуществляться с использованием специальной программы --- чекера.
Каждая задача имеет ограничение по времени и по памяти для выполнения решения,
и система должна корректно обрабатывать случаи превышения времени выполнения или используемой памяти во время выполнения.

\subsubsection*{История дуэлей}
Система должна предоставлять возможность просмотра истории прошедших дуэлей.
Пользователь должен иметь возможность просматривать посылки, сделанные во время участия в дуэли,
а также отправлять решения на проверку после завершения дуэли.

\subsubsection*{Анализ поведения участников}
Система должна собирать информацию о действиях пользователя во время участия в дуэли,
включая события редактирования кода и отправки решений.
На основе собранной информации система должна иметь возможность анализировать поведение участников
и определять степень подозрительности пользователя в конкретной дуэли.

\subsection{Нефункциональные требования}

\subsubsection*{Производительность}

Система должна поддерживать одновременное тестирование нескольких решений.
Результаты тестирования должны отображаться пользователю с минимальной задержкой.

\subsubsection*{Масштабируемость}

Архитектура системы должна поддерживать горизонтальное масштабирование отдельных компонентов при росте нагрузки.

\subsubsection*{Отказоустойчивость}

Система должна быть устойчива к отказу отдельных внутренних и инфраструктурных компонентов.
Система должна обеспечивать возможность восстановления работы после сбоев отдельных компонентов.

\subsubsection*{Безопасность}

Система должна обеспечивать безопасное выполнение пользовательского кода в изолированной среде без доступа к ресурсам хостовой системы.

\subsubsection*{Расширяемость}

Архитектура системы должна поддерживать добавление новых форматов задач без значительного изменения существующих компонентов.
Система должна поддерживать расширение правил проведения дуэлей и турниров.

\subsubsection*{Наблюдаемость}

Система должна поддерживать сбор метрик, журналирование событий и мониторинг состояния компонентов,
а также предоставлять необходимые данные для выявления сбоев.

\subsubsection*{Модульность}

Компоненты системы должны быть слабо связаны между собой и поддерживать независимое развёртывание.